%% Task title and overview, including the significance and what makes it challenging
%% Describe the problem, it’s impact on the real world, and why it is challenging

\section{Overview}

This project focuses on segmenting adult gliomas from MRI scans using a 2D
U-Net convolutional neural network. The BraTS 2021 dataset provides four
co-registered MRI sequences for each subject: T1, T1CE, T2, and T2-FLAIR,
together with an expert-annotated 3D tumor segmentation mask. Our model will
operate on 2D axial slices extracted from these volumes. For each slice, the
corresponding images from all four MRI sequences are stacked as input channels,
and the model predicts which image regions correspond to background or to one
of the tumor sub-regions defined by BraTS.

Manual tumor segmentation requires trained raters to delineate tumor regions
throughout a 3D MRI volume and is therefore time-consuming. An automated model
can generate an initial segmentation that can be reviewed and corrected by an
expert, reducing the amount of manual delineation required. Automated
segmentation is also useful for quantitative research because a fixed model
applies the same learned segmentation procedure to every scan, providing more
repeatable measurements across large datasets than repeated manual
segmentation. This can make it easier to compare tumor regions across subjects
or across scans collected at different time points.

Brain tumor segmentation is challenging because gliomas vary substantially
between subjects in size, shape, location, and appearance. Tumor tissue can
also resemble surrounding healthy tissue, and tumor regions may occupy only a
small fraction of an MRI slice, creating a strong imbalance between tumor and
background. In addition, different tumor sub-regions are distinguished by
patterns across T1, T1CE, T2, and T2-FLAIR rather than by a single MRI sequence
alone. The model must therefore combine information from the four aligned
inputs while identifying irregular tumor boundaries and separating tissue
classes with similar appearances.
